%% ============================================================
%%  CHAPTER 13 — DISCUSSION
%% ============================================================
\chapter{Discussion}
\label{ch:discussion}

The preceding chapters presented this project's methodology and
results in the order they were built: data collection, signal
processing, feature engineering, classification architecture,
results, the control application, and real-time deployment. This
chapter steps back from that build order to interpret what those
results mean, to check the finished system against the research
gaps and contributions claimed at the outset (Chapters~1--2), and
to state plainly where the system falls short. Section~\ref{sec:disc_interpretation}
interprets the headline result in context. Section~\ref{sec:disc_gaps}
revisits the research gaps identified in Chapter~2 against what was
actually built. Section~\ref{sec:disc_metalesson} draws out the
project's central, generalisable engineering lesson.
Section~\ref{sec:disc_evaluation} discusses rigorous evaluation as a
contribution in its own right. Section~\ref{sec:disc_limitations}
presents the project's limitations honestly and in full.
Section~\ref{sec:disc_summary} summarises the chapter.

%% ----------------------------------------------------------
\section{Interpreting the Final Result}
\label{sec:disc_interpretation}
%% ----------------------------------------------------------

The final system reaches 91.3\% accuracy (unweighted \acrshort{loso}
mean) or 91.8\% (pooled), macro-F1 0.913--0.918, on a 5-class
voluntary-artifact command set, evaluated under the strictest
protocol available for a single-subject, cross-day \acrshort{bci}
(Chapters~9--10). Table~\ref{tab:mi_comparison} in Chapter~2 places
this alongside representative motor-imagery \acrshort{bci} systems
from the literature, which report 70--85\% accuracy on 2--4-class
problems using comparable or more sophisticated classifiers
(\acrshort{csp}+\acrshort{lda}, deep CNNs, EEGNet). This project's
result is higher, on more classes, using a comparatively simple
per-branch classifier (nearest-neighbour \acrshort{dtw}, linear
\acrshort{svc} in tangent space) --- but the comparison should be
read carefully rather than as a claim of a generally superior
method.

The reason artifact-based \acrshort{bci} can outperform motor
imagery here is architectural to the \emph{problem}, not just the
\emph{solution}: voluntary blinks and jaw contractions are
high-amplitude, largely voluntary, and physiologically distinct
from one another (\acrshort{eog} vs.\ \acrshort{emg}, Chapter~3),
whereas motor imagery relies on subtle, often unreliable
sensorimotor rhythm modulation that many users struggle to produce
consistently even after extensive training (Chapter~2). This
project's accuracy advantage is therefore best read as evidence
that \textbf{trading generality (imagined movement of any limb) for
specificity (five deliberately chosen, physiologically distinct
gestures)} is a favourable trade for a prosthetic-control use case,
not as evidence that the classification methods developed here
would transfer that same advantage to a motor-imagery problem.
Whether that trade is acceptable depends entirely on the deployment
context: a user with reliable voluntary control over eye and jaw
muscles is a good match for this system; a user without such
control is not, and would need a motor-imagery or hybrid approach
instead.

%% ----------------------------------------------------------
\section{Revisiting the Research Gaps}
\label{sec:disc_gaps}
%% ----------------------------------------------------------

Chapter~2 identified six research gaps (G1--G6) in the
artifact-based \acrshort{bci} literature and mapped this project's
seven claimed contributions (C1--C7, Chapter~1) onto them.
Table~\ref{tab:gaps_delivered} restates that mapping alongside the
specific chapter and result that backs each claim, so that every
contribution claimed in Chapter~1 can be checked against concrete
evidence rather than taken on faith.

\begin{table}[htbp]
\centering
\caption{Research gaps (Chapter~2) and contributions (Chapter~1),
cross-referenced against the concrete evidence produced for each.}
\label{tab:gaps_delivered}
\renewcommand{\arraystretch}{1.35}
\begin{tabular}{c p{3.6cm} p{3.0cm} p{3.0cm}}
\toprule
\textbf{Gap} & \textbf{Claim} & \textbf{Evidence} &
\textbf{Chapter} \\
\midrule
G1 & Physiology-matched hybrid architecture (C1) & +4.1\% Riemann,
    +2.1\% \acrshort{dtw}, over flat/hierarchical baseline & 9, 10 \\
G2 & Rest-window contamination addressed (C4) & Activity gate,
    +3.1\% accuracy & 7, 10 \\
G3 & Data leakage in evaluation addressed (C5) & 4-test audit,
    shuffle test at 21\%, \acrshort{loso} adopted & 9 \\
G4 & Novel L/R bruxism asymmetry feature (C2) & T3--T4 index,
    bruxism\_left F1 +0.128 at introduction & 8 \\
G5 & End-to-end system with \acrshort{fsm} + activity gate (C6) &
    Full pipeline, 14\,ms latency, 3-consecutive confirmation &
    11, 12 \\
G6 & Consumer-hardware validation (C7, partial) & 19-channel
    device; Emotiv target hardware not yet validated (see
    Section~\ref{sec:disc_limitations}) & 4, 6 \\
\bottomrule
\end{tabular}
\end{table}

Five of the six gaps (G1--G5) are addressed with direct, quantified
evidence produced over the course of this project. G6 is only
\textbf{partially} addressed, and this project does not claim
otherwise: the 19-channel medical-grade device used for data
collection (Chapter~6) is a mid-range research device, not the
lower-channel-count consumer hardware (the Emotiv EPOC~X) that was
always the intended deployment target. This gap between
data-collection hardware and deployment hardware is discussed fully
as a limitation in Section~\ref{sec:disc_limitations}, since it is
the single largest piece of unfinished validation work remaining in
the project.

%% ----------------------------------------------------------
\section{The Central Lesson: Architecture Over Features}
\label{sec:disc_metalesson}
%% ----------------------------------------------------------

If this project has one methodological lesson worth generalising
beyond its specific gesture set and hardware, it is the one first
surfaced in Chapter~8 and confirmed twice more in Chapter~9:
\textbf{when a multi-class problem contains classes with
fundamentally different underlying signal physiology, no amount of
additional feature engineering on a single flat classifier can
resolve the resulting tension, because a flat classifier has no
mechanism to apply a feature conditionally.}

This pattern appeared three separate times in this project, each
time independently:

\begin{itemize}[leftmargin=*]
    \item \textbf{Feature engineering (Chapter~8):} every advanced
    feature family tested (\acrshort{dwt}, \acrshort{tkeo},
    \acrshort{emg} descriptors) helped one class family and hurt
    the other, never both --- because muscle-oriented features are
    meaningless for \acrshort{eog} events and vice versa.
    \item \textbf{Muscle classification (Chapter~9):} amplitude-based
    features (\acrshort{tkeo}) failed to survive cross-session
    drift; only when the \emph{representation itself} was changed
    to something inherently insensitive to amplitude (Riemannian
    covariance geometry) did the muscle branch improve.
    \item \textbf{Blink classification (Chapter~9):} four
    independently-designed counting-based methods all failed on
    contaminated data for the same underlying reason; only when the
    \emph{representation itself} was changed from counting discrete
    events to comparing waveform shape (\acrshort{dtw}) did the
    blink branch improve.
\end{itemize}

In each case, the fix that worked was not ``try harder'' within the
existing representation --- more features, more hyperparameter
tuning, a bigger classifier --- but a change of representation
matched to the physiology of the specific signal being classified.
This is the throughline connecting Chapters~8 and~9, and it is
offered here as a lesson that should generalise to other multi-class
\acrshort{bci} and biosignal problems with heterogeneous class
physiology, not just to this project's specific five gestures.

%% ----------------------------------------------------------
\section{Rigorous Evaluation as a Contribution in Itself}
\label{sec:disc_evaluation}
%% ----------------------------------------------------------

The evaluation-methodology crisis documented in Chapter~9 --- an
apparently strong $\sim$80--85\% result that turned out to be
substantially inflated by epoch-overlap leakage --- is presented in
this report not as an embarrassing false start to be minimised, but
as one of the project's more transferable contributions (C5,
Chapter~1). Two aspects of the audit are worth drawing out
explicitly here:

\begin{enumerate}[leftmargin=*]
    \item \textbf{The shuffle test is cheap and diagnostic.}
    Permuting class labels while leaving every other part of the
    pipeline untouched, and confirming the result collapses to
    chance (21\% for five classes, Chapter~9), is a small addition
    to any evaluation pipeline that provides strong evidence the
    pipeline is not manufacturing signal through some
    implementation artefact. This is a generally applicable sanity
    check, independent of the specific classes or features involved,
    and is recommended here as a default step in any \acrshort{bci}
    evaluation pipeline rather than a project-specific one-off.
    \item \textbf{The right level of strictness depends on the
    deployment question, not on what is easiest to report.} This
    project deliberately reports \acrshort{loso} --- the strictest
    protocol available given three sessions --- as its headline
    number, rather than a less strict but higher-scoring protocol
    (Chapter~9, Table~\ref{tab:honest_protocols}). The gap
    between GroupKFold's development-time 88.0\% and \acrshort{loso}'s
    91.3\% is a reminder that even a leak-free protocol can
    understate or overstate deployment performance if it does not
    match the actual deployment question (``how well does this work
    on a new day for the same user?'').
\end{enumerate}

Chapter~2 identified underreporting of evaluation leakage as a gap
in the wider artifact-based \acrshort{bci} literature (G3). This
project's contribution to that gap is not merely ``we did not leak''
but the fully worked audit methodology itself --- the four
diagnostic tests, the shuffle test, and the explicit protocol
comparison table --- which is offered as a template other
\acrshort{bci} projects using overlapping sliding-window epochs
could adopt directly.

%% ----------------------------------------------------------
\section{Limitations}
\label{sec:disc_limitations}
%% ----------------------------------------------------------

This report has aimed throughout to document failures and
unresolved issues alongside successes (Chapters~6--12); this section
consolidates the limitations that remain unresolved at the time of
writing, rather than scattering them across earlier chapters.

\subsection{Single-Subject Design}

Every result in this report describes a model calibrated to, and
evaluated on, one individual (Chapter~6). This is standard practice
for \acrshort{bci} systems intended to be personalised per user
(Chapter~2), and \acrshort{loso}'s cross-\emph{session} evaluation
is the correct generalisation target for that design --- but it
means no claim in this report should be read as evidence of
cross-\emph{subject} accuracy. A new user would need to collect
their own multi-session calibration dataset and retrain before this
system's accuracy figures would apply to them.

\subsection{Residual Single-Blink Contamination}

\acrshort{dtw} shape-matching (Chapter~9) resolved the
\emph{architectural} ceiling that four counting-based methods could
not get past, but it did not retroactively clean the underlying
recordings: roughly 45\% of \code{single\_blink} recordings still
physically contain multi-blink clusters (Chapter~6), and the
confusion matrix (Chapter~10) shows this residual contamination is
still the dominant source of error in the final model (82 of 102
total misclassifications). The architecture-level ceiling has been
resolved; the underlying data-quality problem has not, and the
single highest-leverage remaining improvement to this system's
accuracy would likely be a clean re-recording of the
\code{single\_blink} class with deliberate, verified pauses between
individual blinks (Chapter~14), rather than any further
classifier-side change.

\subsection{Data-Collection vs.\ Deployment Hardware Mismatch}

The complete dataset (Chapter~6) was recorded on a 19-channel
research-grade device sampling at \SI{200}{\hertz}. The intended
deployment target is the Emotiv EPOC~X, a 14-channel consumer
headset with a different electrode layout and an on-device 45\,Hz
hardware filter. Channel correspondence between the two devices is
only partial: T3/T4 (used throughout this report for jaw
\acrshort{emg}) map approximately to the Emotiv's T7/T8, and Fp1/Fp2
(used for blink \acrshort{eog}) map approximately to AF3/AF4, but
neither mapping is exact, and no data has yet been collected on the
target device to validate that the trained model transfers. This is
the single largest piece of unfinished validation in the project:
every accuracy figure in this report describes performance on the
data-collection hardware, not the deployment hardware, and the two
should not be assumed equivalent without direct verification
(Chapter~14).

\subsection{Incomplete Firmware and Application Integration}

As documented in full and without qualification in Chapter~11, only
the firmware scaffold (Layer~0) is verified running on the physical
ESP32 at the time of writing. Servo control, the WiFi access point,
and the live mapping/pose editors --- the layers that would let a
user actually configure and observe the system through the web app
described in Chapter~11 --- are designed but not yet built. The web
application itself is a complete, working prototype, but has only
been exercised in its mock-data mode, not against the real firmware
or a live classifier stream. The real-time inference pipeline
(Chapter~12) and the firmware (Chapter~11) have not yet been run
together end-to-end on physical hardware; each has been validated
independently (inference against the offline-trained model and
recorded data; firmware Layer~0 on the physical board) but not yet
as a connected whole.

\subsection{Real-Time Activity Gate Fragility}

The real-time activity gate (Chapter~12) uses a fixed
20\,\si{\micro\volt} threshold, calibrated against the amplitude
distributions observed in the recorded dataset. Unlike the
recording-relative offline gate (Chapter~7), this fixed threshold
cannot adapt to a new session's specific electrode impedance or
amplifier gain, and its continued validity in a live deployment
session has not been verified outside the conditions under which it
was calibrated.

%% ----------------------------------------------------------
\section{Chapter Summary}
\label{sec:disc_summary}
%% ----------------------------------------------------------

This chapter interpreted the project's 91.3--91.8\% \acrshort{loso}
result in the context of the motor-imagery \acrshort{bci}
literature (Chapter~2), attributing its accuracy advantage to the
deliberate choice of high-amplitude, physiologically distinct
voluntary gestures rather than to a generally superior
classification method. Five of the six research gaps identified in
Chapter~2 were shown to be addressed with direct, quantified
evidence; the sixth (consumer-hardware validation) was shown to be
only partially addressed, pending validation on the intended Emotiv
EPOC~X deployment hardware. The project's central, generalisable
engineering lesson --- that heterogeneous-physiology multi-class
problems are better solved by routing to a physiology-matched
representation than by adding features to a flat classifier --- was
drawn out explicitly across its three independent appearances in
Chapters~8 and~9. The evaluation-methodology audit (Chapter~9) was
presented as a contribution in its own right, with the shuffle test
recommended as a generally applicable diagnostic beyond this
project. Finally, five specific, unresolved limitations were stated
plainly: the single-subject design, residual single-blink
contamination not fixable by further classifier changes, the
data-collection/deployment hardware mismatch, incomplete firmware
and application integration, and the real-time activity gate's
fixed-threshold fragility. The following, final chapter concludes
the report and lays out the concrete future work implied by each of
these limitations.